\documentclass[11pt]{article}
\usepackage[margin=1in]{geometry}
\usepackage{mathpazo}
\usepackage{amsmath,amssymb}
\usepackage{microtype}
\usepackage{booktabs}
\usepackage{graphicx}
\usepackage[dvipsnames]{xcolor}
\usepackage{tikz}
\usetikzlibrary{positioning,arrows.meta,calc,fit,backgrounds,decorations.pathreplacing}
\usepackage[most]{tcolorbox}
\usepackage{enumitem}
\usepackage{listings}
\usepackage[font=small,labelfont=bf]{caption}
\usepackage{hyperref}
\hypersetup{colorlinks=true,linkcolor=MidnightBlue,citecolor=MidnightBlue,urlcolor=MidnightBlue}
\usepackage{cleveref}
\urlstyle{tt}
\def\UrlBreaks{\do\_\do/\do.}
\newcommand{\file}[1]{\nolinkurl{#1}}
\newcommand{\code}[1]{\texttt{#1}}
\newcommand{\mhh}{m_{hh}}
\newcommand{\bC}{\boldsymbol{C}}
\newcommand{\snoise}{s_{\text{noise}}}

\definecolor{trainable}{HTML}{1B6CA8}
\definecolor{physics}{HTML}{2E933C}
\definecolor{datac}{HTML}{7B4EA3}
\definecolor{warn}{HTML}{C0392B}

\newtcolorbox{keyidea}[1][]{enhanced,breakable,colback=physics!6,colframe=physics!70!black,
  fonttitle=\bfseries,title={#1},boxrule=0.6pt,arc=2pt,left=6pt,right=6pt}
\newtcolorbox{aside}[1][]{enhanced,breakable,colback=trainable!5,colframe=trainable!70!black,
  fonttitle=\bfseries,title={#1},boxrule=0.6pt,arc=2pt,left=6pt,right=6pt}
\newtcolorbox{assume}[1][]{enhanced,breakable,colback=datac!5,colframe=datac!70!black,
  fonttitle=\bfseries,title={#1},boxrule=0.6pt,arc=2pt,left=6pt,right=6pt}
\newtcolorbox{watch}[1][]{enhanced,breakable,colback=warn!5,colframe=warn!70!black,
  fonttitle=\bfseries,title={#1},boxrule=0.6pt,arc=2pt,left=6pt,right=6pt}
\setlist{itemsep=2pt,topsep=3pt}
\setlength{\emergencystretch}{3em}

\lstset{language=Python, basicstyle=\ttfamily\small, columns=fullflexible,
  keepspaces=true, commentstyle=\color{physics!70!black}\itshape,
  keywordstyle=\color{trainable!90!black}\bfseries, stringstyle=\color{datac},
  showstringspaces=false, frame=single, rulecolor=\color{gray!40},
  backgroundcolor=\color{gray!3}, xleftmargin=4pt, xrightmargin=4pt,
  aboveskip=6pt, belowskip=6pt}

\title{\textbf{Evaluating the density network}\\[4pt]
\large Plan: the same tests, the same measures, the same table as the smooth model}
\author{P.~Bhat (Fermilab), N.~De~Filippis (Politecnico di Bari),\\
K.F.~Johnson, H.B.~Prosper, R.~Starnes (Florida State University)}
\date{5 November 2026}

\begin{document}
\maketitle

\begin{abstract}
\noindent
The fine-tuned density networks (\file{runs/FTD_SiLU_*_pearson}, started from
\file{runs/DENS_L1_SiLU_thr}) look good in the fine-tune notebook's own measure. That measure
is Neyman's $\chi^2$ with an estimated number of events, which is not what the paper uses. We
propose to evaluate them exactly as the smooth model was evaluated: the likelihood-ratio
statistic with the exact Poisson counts, calibrated pulls with a model-error term determined on
data not used in the fit, and a direct comparison of the densities. The work goes into one new
notebook, \file{code/heft_density_eval.ipynb}, and a short results note. Decisions for you are in
\cref{sec:decisions}.
\end{abstract}

\tableofcontents

% ===========================================================================
\section{Why a new evaluation}

\begin{aside}[Two measures of agreement with a count]
Let $n$ be the POWHEG count in a bin and $\mu$ the model's prediction, in events.
\begin{itemize}
  \item \textbf{Neyman's $\chi^2$}, $(n-\mu)^2/n$, uses the \emph{observed} count as the variance.
    The fine-tune notebook uses this, with $n$ estimated through $N_\text{eff}$. It favours
    predictions that are biased low, because a downward fluctuation gets a smaller variance and
    so a larger weight.
  \item \textbf{The likelihood-ratio statistic}, $-2\ln\lambda = 2[\mu - n + n\ln(n/\mu)]$, the
    logarithm of $\mathrm{Poisson}(n;\mu)/\mathrm{Poisson}(n;n)$, is the exact measure for counts.
    For a correct model it follows $\chi^2_1$ (mean 1). The paper uses this, with the exact counts
    recovered by \file{dihiggs/counts.py}.
\end{itemize}
A model fitted with its own Poisson variances (\code{VARIANCE\_FROM='model'}), as the density
networks were, is penalised by Neyman's measure for \emph{not} being biased low. The comparison
with the smooth model must therefore be redone with $-2\ln\lambda$.
\end{aside}

The reference numbers on the 1,000 test spectra (mean $-2\ln\lambda$ per bin):

\begin{center}
\small
\begin{tabular}{lccc}
\toprule
model & 255--555 GeV & $>555$ GeV & median per spectrum\\
\midrule
smooth model & 1.566 & 1.100 & 1.078\\
free fit per bin (the floor for any 23-term model) & 1.559 & 1.105 & 1.075\\
old network ensemble (20 networks) & 1.561 & 1.121 & 1.085\\
\bottomrule
\end{tabular}
\end{center}

% ===========================================================================
\section{The tests}
\label{sec:tests}

All on the 1,000 test spectra, which played no part in training, for the L1 starting network and
the four fine-tunes ($\delta = 1, 2, 3$ and plain $\chi^2$).

\subsection{Accuracy: $-2\ln\lambda$}

Mean per bin below and above 555~GeV, and the median per spectrum, as in the table above; and per
bin as a function of $\mhh$ (the figure \file{fig_app_accuracy} with the density network added).
The density network's prediction for a bin is its bin integral, \code{HEFTDensity.spectrum(theta)}.

\subsection{Calibration: a model-error term for the network}

The pulls need the total uncertainty of a prediction. For the smooth model this is the Poisson
noise, the exact fit uncertainty, and the model-error term $\delta = \varepsilon(m)\sum_i|C_iA_i|$
(\file{docs/dihiggs_model.pdf}, section 6).

\begin{itemize}
  \item \textbf{Fit uncertainty.} A single network has none. The ensembles showed that the network
    spread is a few per cent of the noise (5.6\%), so we set it to zero for now and note this. A
    seeds ensemble (\cref{sec:later}) will supply it.
  \item \textbf{Model-error term.} We determine the network's own $\varepsilon(m)$ exactly as for
    the smooth model: on the 38,076 quick-test rows, never used in training, solve
    $\langle (t-f)^2/(\snoise^2 + \varepsilon_k^2 S^2)\rangle = 1$ in each bin, with
    $S = \sum_i |C_i A_{ik}|$ from the network's own $A_{ik}$, then smooth $\ln\varepsilon$ by a
    quadratic. As a check we also use the smooth model's $\varepsilon$ unchanged.
\end{itemize}
Measures: standard deviation of the pull overall and below and above 555~GeV; the fraction with
$|z|<2$; $\langle z^2\rangle$ by fifth of the cancellation ratio. The network's $\varepsilon(m)$ is
itself a result: if the network sits at the same floor as the smooth model, its $\varepsilon$ should
be similar.

\subsection{The density itself}

\begin{itemize}
  \item \textbf{Against the smooth model}, on a 0.5~GeV grid from 255 to 1455~GeV, for all 1,000 test
    points: the relative difference of $d\sigma/dm$, at the bin centres and between them, and the
    fraction of points within the smooth model's total $\pm1\sigma$ and $\pm2\sigma$ bands.
  \item \textbf{Between bin centres}: the figure \file{fig_interpolation} redrawn with the density
    network beside the old ensemble and the smooth model, for the same three spectra.
  \item \textbf{The 23 coefficient functions} $a_i(m)$ of the network, overlaid on the smooth model's
    bands (\file{fig_coefficients}). The poorly determined $a_{16}$--$a_{20}$ may differ; the
    well-determined ones should agree.
\end{itemize}

\subsection{Below the fitted range}

With the threshold factor, the density network is defined down to 250~GeV. Its integral over
250--255~GeV is a prediction for the unused POWHEG bin 240--255~GeV (the source file holds it). The
smooth model predicted that bin 8\% high in total over the test spectra; we compare.

% ===========================================================================
\section{Implementation}

\begin{itemize}
  \item \textbf{New notebook \file{code/heft_density_eval.ipynb}}, in the style of the others (header,
    banner comments, switches). It reuses \file{docs/paper/smooth_common.py} for the test spectra with
    their exact counts, the quick-test rows, the smooth model's predictions and the old ensemble.
    A new helper there, \code{load\_density\_network(run)}, rebuilds a \code{HEFTDensity} from a
    run's \file{dihiggs_config.yaml} (activation, threshold, nodes, input) and loads its parameters.
  \item \textbf{Switches}: \code{RUNS} (the run folders to evaluate), \code{FEATURED} (the run shown in
    the figures, default \file{FTD_SiLU_huber_d2_pearson}).
  \item \textbf{Outputs}: printed tables; figures in \file{docs/density/} (accuracy versus $\mhh$,
    calibration by cancellation fifth, density comparison, coefficients, interpolation).
  \item \textbf{Results note} \file{docs/density/density_results.pdf}, in the house style.
  \item Runs on a laptop CPU in a few minutes; no GPU needed.
\end{itemize}

% ===========================================================================
\section{Later, on the 5090}
\label{sec:later}

\begin{enumerate}
  \item A seeds ensemble of the featured configuration (10 seeds): the network's fit uncertainty, and
    its spread between bin centres on real data.
  \item The merged-bin test: train on 30~GeV bins, predict the 15~GeV bins (the density model only
    needs 30~GeV edges; the data preparation merges pairs of rows, whose counts add).
\end{enumerate}

% ===========================================================================
\section{Decisions for you}
\label{sec:decisions}

\begin{keyidea}[Decisions before I proceed]
\begin{enumerate}
  \item \textbf{Evaluate all five runs}, featuring $\delta = 2$ in the figures. Agreed, or feature
    another?
  \item \textbf{Network fit uncertainty set to zero} until a seeds ensemble exists. Agreed?
  \item \textbf{The network's own $\varepsilon(m)$} from the quick-test rows, with the smooth model's
    $\varepsilon$ as a check. Agreed?
  \item \textbf{The 240--255~GeV check} included. Agreed?
  \item \textbf{A new notebook} \file{code/heft_density_eval.ipynb} and a helper in
    \file{smooth_common.py}, rather than additions to \file{heft_test.ipynb}. Acceptable?
\end{enumerate}
\end{keyidea}

\end{document}
